% Pizarra de T3, páginas 46--47.
\begin{ejer}[El Corolario central: la equivalencia contravariante]\label{ejer:t3-equivalencia-afines}
Muestra que la categoría de anillos y la categoría opuesta de esquemas afines son equivalentes. \hfill [Vakil 7.3.E]
\end{ejer}
\begin{sol}
Probemos que $\Spec\colon\cat{Ring}^{\mathrm{op}}\to\cat{LocAn}$, definido en el \cref{ejer:t3-spec-funtor}, es plenamente fiel. Veamos que $(\pi,\pi^\#)\mapsto\pi^\#_{\Spec B}$ es inverso de
\[
\Hom(B,A)\longrightarrow\Hom(\Spec A,\Spec B).
\]
\begin{itemize}
\item Para una de las composiciones, sea $(\pi,\pi^\#)\colon\Spec A\to\Spec B$. Entonces la flecha
\[
\varphi\coloneq\pi^\#_{\Spec B}\colon
B=\mathcal O_{\Spec B}(\Spec B)\longrightarrow\mathcal O_{\Spec A}(\Spec A)=A
\]
cumple, para todo $\mathfrak p\en\Spec A$,
\[
\begin{aligned}
a\en\varphi^{-1}(\mathfrak p)
&\quad\Longleftrightarrow\quad\frac{\varphi(a)}{1}\en\mathfrak p A_{\mathfrak p}
\quad\Longleftrightarrow\quad\pi^\#_{\mathfrak p}\left(\frac{a}{1}\right)\en\mathfrak p A_{\mathfrak p}\\
&\quad\Longleftrightarrow\quad\frac{a}{1}\en\pi(\mathfrak p)B_{\pi(\mathfrak p)}
\quad\Longleftrightarrow\quad a\en\pi(\mathfrak p).
\end{aligned}
\]
Aquí usamos, respectivamente, las definiciones del mapa en tallos y la localidad de $\pi^\#_{\mathfrak p}$. Se sigue $\pi(\mathfrak p)=\varphi^{-1}(\mathfrak p)$, así $\pi=\Spec\varphi$ en espacios.

Para gavillas, sea $g\en B$ arbitrario; el siguiente diagrama conmuta por naturalidad de $\pi^\#$:
\[
\begin{tikzcd}[column sep=large,row sep=large]
B\arrow[r,"{\varphi=\pi^\#_{\Spec B}}"]\arrow[d]&A\arrow[d]\\
B_g\arrow[r,"\pi^\#_{D(g)}"']&A_{\varphi(g)}.
\end{tikzcd}
\]
Esto implica que $\pi^\#_{D(g)}$ es extensión de $B\xrightarrow{\varphi}A\to A_{\varphi(g)}$ para todo abierto distinguido $D(g)$, por lo que $\pi^\#=\varphi^\#$, donde $\varphi^\#$ denota el mapa de gavillas construido en el \cref{ejer:t3-spec-funtor}. Hemos probado
\[
(\Spec\varphi,\varphi^\#)=(\pi,\pi^\#).
\]

\item Para la otra composición, sea $\pi^\#\colon B\to A$ morfismo de anillos. Este induce un
\[
(\pi,\pi^\#)\colon\Spec A\longrightarrow\Spec B.
\]
Entonces basta ver que el mapa inicial $\pi^\#$ coincide con $\pi^\#_{\Spec B}$. En efecto, para todo $b\en B$,
\[
\pi^\#_{\Spec B}\left(\frac{b}{1}\right)
=\frac{\pi^\#(b)}{\pi^\#(1)}=\frac{\pi^\#(b)}{1},
\]
recordando que estamos identificando $A=A_1$ y $B=B_1$. Esto prueba que $\pi^\#=\pi^\#_{\Spec B}$.
\end{itemize}
Finalmente, tenemos que $\Spec$ induce una biyección
\[
\Hom(B,A)\xlongrightarrow{\cong}\Hom(\Spec A,\Spec B),
\]
por lo que es pleno y fiel. Como su imagen esencial es, por definición, la categoría de esquemas afines, tenemos un funtor
\[
\Spec\colon\cat{Ring}^{\mathrm{op}}\longrightarrow\cat{AffSch}
\]
plenamente fiel y esencialmente sobre: una equivalencia de categorías. Más aún, la asignación
\[
(\pi,\pi^\#)\longmapsto\pi^\#_{\Spec B}
\]
es la acción en flechas del funtor $\Gamma$ de secciones globales, por lo que hemos probado que $\Gamma\colon\cat{AffSch}\to\cat{Ring}^{\mathrm{op}}$ es un cuasi-inverso de $\Spec$.
\end{sol}
